单个协议输入所满足的安全性质，并不能单独确定多个输入以何种方式组合到一次协议操作中仍然安全。当接收方把若干与发送方相关联的输入归入同一处理过程时，即使每个条目都有合法来源，整个集合仍可能重复出现接口原本只允许一次的参与方。安全消息系统中的会话处理研究表明，组合边界需要与其组成元素分别考察\cite{cremers2023session}。本文关注的并非跨会话或跨会话群组的组合，而是一次接收方批处理内部的输入关系。

\kwaay 是一种后量子、类似X3DH且具有可否认性的认证密钥交换协议\cite{collins2024kwaay}。其\batchreceive 操作在接收方状态下成组处理多个输入；完整构造明确规定，同一次调用中的各元素应对应不同参与方\cite[第5.1节，第21页]{collins2024kwaayfull}。这一条件已经存在于源协议。本文所考察的是它维护了何种关系，以及作用于其他坐标的限制能否维护同一关系。下文将该批内关系记为\distinctparty；这一分析性名称由本文引入，并不表示K-Waay遗漏了新的协议条件。

区别这些坐标具有直接意义。同一建模参与方可以在不同发送实例中产生两条消息。由这两个条目组成的批次可以同时满足消息不同、任一精确来源均未被重复接受，却仍在参与方层面出现两次。反过来，把同一个精确条目放入两个槽位，则引出的是接受多重性问题。因而，在判断某项接纳限制能否达到目标之前，必须区分参与方身份、消息身份以及发送实例（occurrence）匹配。身份概念与单射对应关系已有系统研究\cite{lupetti2006names,lowe1997hierarchy}，本文关注的是这些关系在K-Waay批处理接口上的相互作用。

\paragraph{研究问题与方法。}
围绕上述边界，第\ref{sec:identity-problem}节提出三个问题：放宽不同参与方条件后，哪些参与方重复组合变为可达（RQ1）；在该抽象下，一个匹配的发送来源能否支持同一批次中的多个接收事件（RQ2）；仅约束消息坐标是否足以推出批内参与方区分（RQ3）。分析比较三个固定两槽Tamarin模型\cite{meier2013tamarin}：宽松接纳模型、消息级限制模型和参与方级限制模型。三者共享发送生成、输入收集、精确来源匹配及顺序处理生命周期，仅在接纳谓词及相应拒绝分支上约束不同坐标。参与方级模型充当正控制，用于直接检验目标关系，并不预设直接比较参与方身份是维护该关系的唯一实现方式。

\paragraph{结果与贡献。}
本文的贡献在于形式化阐明原协议所述条件的语义作用，并在明确的建模假设下分离消息级、发送实例级和参与方级性质。宽松接纳模型存在一条执行：一个匹配的\textsf{Send}来源在同一批次上下文中支持两个\receiveraccept 事件。消息级模型能够保持消息区分和作用域内的精确来源单射性，却仍可共同接纳来自同一参与方的两条不同消息。参与方级模型直接约束所收集条目的参与方坐标；已有结果既确认接收事件之间的参与方区分，也给出可达的有效异参与方批次。其中，消息级结果给出了核心分离：消息区分和作用域内的单射性成立，并不蕴含批内参与方区分。

\paragraph{研究边界。}
分析假定对手在建模边界上能够控制候选条目的批次组合，但不据此认定网络攻击者能够控制真实部署中的批次构造器。新鲜符号消息和发送实例坐标、理想化的持久来源关系以及模型局部的接收事件，使三个模型之间的比较保持简洁、可解释。当前工作给出源接口到抽象对象的映射，但没有建立完整K-Waay执行到模型的精化定理；其结果既不是密码学攻击或部署漏洞，也不是针对任意批大小的机器证明。论文使用既有模型执行记录，2026年9月16日完成的独立复跑又保存了精确命令、输入哈希、证明器标准输出与错误输出以及导出轨迹；该复跑验证了当前抽象结果的可复现性，但不恢复历史运行的出处，也不扩展研究结论。

第\ref{sec:identity-problem}节介绍协议接口及身份区分问题，第\ref{sec:objective-threat-model}节定义分析目标和对手边界。后续第3、4节分别给出形式化抽象与机器分析，第5、6节讨论设计含义、局限和相关工作，第7节总结全文。
